% !PRINTED 344
\weberpage{344}
%<<CONT
marked, an approximation to the absolutely smallest root. But by shifting the origin one can make any root the absolutely smallest, though for this, admittedly, the knowledge of a value approximated to a certain degree is necessary.

A formula that suffices for all cases has been given by Fr. \gesperrt{Meyer}\footnote{Mathematische Annalen, Bd. 33.}. Let $\alpha_1, \alpha_2, \alpha_3 \ldots \alpha_n$ be the roots of the equation $f(x) = 0$, which may be real or imaginary. One seeks a point $p$ in the $x$-plane such that a circle can be described about $p$ as centre which contains \gesperrt{only one} of the points $\alpha_1, \alpha_2 \ldots \alpha_n$, say the point $\alpha_1$, so that the absolute values of
\weq{5}{
\frac{p-\alpha_1}{p-\alpha_2}, \quad \frac{p-\alpha_1}{p-\alpha_3}, \quad \cdots \quad \frac{p-\alpha_1}{p-\alpha_n}
}
are proper fractions. Such points always exist; if need be, they must be found by the \gesperrt{Sturm} method. If, in addition, one chooses an arbitrary function $\varphi(x)$ that has no common divisor with $f(x)$, e.g. $\varphi(x) = 1$, and forms the symmetric functions
\weq{6}{
y_m = \dfrac{\dfrac{\alpha_1\varphi(\alpha_1)}{(p-\alpha_1)^m} + \dfrac{\alpha_2\varphi(\alpha_2)}{(p-\alpha_2)^m} + \cdots + \dfrac{\alpha_n\varphi(\alpha_n)}{(p-\alpha_n)^m}}{\dfrac{\varphi(\alpha_1)}{(p-\alpha_1)^m} + \dfrac{\varphi(\alpha_2)}{(p-\alpha_2)^m} + \cdots + \dfrac{\varphi(\alpha_n)}{(p-\alpha_n)^m}},
}
then with increasing $m$ these approach the limit $\alpha_1$, and the more rapidly, the smaller the absolute values of the fractions (5) already are.

The correctness of this shows itself at once if one writes the formula (6) in the following way:
\[
y_m = \frac{\alpha_1\varphi(\alpha_1) + \alpha_2\varphi(\alpha_2)\left(\dfrac{p-\alpha_1}{p-\alpha_2}\right)^m + \cdots + \alpha_m\varphi(\alpha_m)\left(\dfrac{p-\alpha_1}{p-\alpha_n}\right)^m}{\varphi(\alpha_1) + \varphi(\alpha_2)\left(\dfrac{p-\alpha_1}{p-\alpha_2}\right)^m + \cdots + \varphi(\alpha_m)\left(\dfrac{p-\alpha_1}{p-\alpha_n}\right)^m}.
\]
% [sic] source prints the final outer subscripts as $m$, although the sequence and the inner $p-\alpha_n$ indicate $n$

\begin{center}
\S.\,111.

\gesperrt{The approximation method of Gräffe.}
\end{center}

On an idea similar to that of the \gesperrt{Bernoulli} approximation there also rests a procedure that \gesperrt{Gräffe}
%CONT>>
